\documentclass[10pt,oneside,notitlepage,a4paper]{book}
\usepackage[no-math]{fontspec}
\setmainfont{Times New Roman}
\setmonofont{Consolas}
\usepackage{amsmath,amsfonts,amssymb,amsthm}
\usepackage[xcharter,nosymbolsc]{newtxmath}
\DeclareMathSizes{11}{11}{9}{7}
\DeclareMathSizes{10}{10}{7}{5}
\DeclareMathSizes{9}{9}{7}{5}
\SetSymbolFont{operators}{normal}{OT1}{cmr}{m}{n}
\SetMathAlphabet{\mathrm}{normal}{OT1}{cmr}{m}{n}
\SetMathAlphabet{\mathbf}{normal}{OT1}{cmr}{m}{n}
\usepackage[hidelinks]{hyperref}
\renewcommand{\thesection}{\arabic{section}}
\renewcommand{\theequation}{\arabic{equation}}
\renewcommand{\bibname}{References}
\pagestyle{plain}
\newenvironment{abstract}
  {\begin{quote}\small\textbf{Abstract.}\ }
  {\end{quote}}
\setlength{\textwidth}{155mm}
\setlength{\textheight}{235mm}
\setlength{\oddsidemargin}{2mm}
\setlength{\evensidemargin}{2mm}
\setlength{\topmargin}{-8mm}
\setlength{\parskip}{0.25em}
\linespread{1.08}
\allowdisplaybreaks[2]
\newtheorem{theorem}{Theorem}[section]
\newtheorem{lemma}[theorem]{Lemma}
\newtheorem{proposition}[theorem]{Proposition}
\theoremstyle{remark}
\newtheorem{remark}[theorem]{Remark}
\newcommand{\N}{\mathrm N}
\newcommand{\Rea}{\operatorname{Re}}
\newcommand{\one}{\mathbf 1}
\newcommand{\eps}{\varepsilon}
\newcommand{\pp}{\mathfrak p}
\newcommand{\OO}{\mathcal O_F}
\newcommand{\st}{\sigma_\tau}
\newcommand{\bp}{\beta_*}
\title{Beyond the Seven-Eighths Barrier:\\[0.2em] A Sharper Quasi-Riemann Boundary}
\author{Argonaut Math}
\date{October 7, 2026}

\begin{document}
\fontsize{11}{15}\selectfont
\maketitle
\fontsize{11}{15}\selectfont

\begin{abstract}
We refine the zero-free boundary in OpenAI's September 30 quasi-Riemann
paper by slightly increasing the total length of the compensation primes
and adjusting the two column scales. From that paper's stated unconditional
results, we obtain the strict zero-free half-plane
$\Rea s>3499999/4000000$ for finite-order Hecke $L$-functions over
$\mathbb Q(\sqrt{-3})$ and for all Dirichlet $L$-functions. This improves
the boundary $7/8$ by $1/4000000$. The arithmetic and analytic estimates still concern the same
finite probe. On the arithmetic side, its rescaling coefficients absorb
the extra completed-row energy. On the analytic side, an explicit margin
in the central-row estimate remains positive. We give the affected
counting, contour, and principal-term arguments in full. The parameters
are chosen so that the moment mesh is fixed independently of the later
principal approximation. The foundational reflection, sieve, and moment
theorems are taken from the original paper; no new short mean-square
assumption is imposed.
\end{abstract}

\section{The boundary and the perturbation}

This paper refines the parameters in OpenAI's September 30 argument
as part of the Argonaut Math research program. That argument proves a
zero-free half-plane with boundary $7/8$. It uses a finite arithmetic
probe with two representations: one admits a sieve estimate, and the other
produces an analytic signal. We show that the construction permits a small
increase in the length of the compensation primes. Adjusting the column
scales to accommodate that increase gives a strictly smaller boundary.

Riemann's 1859 paper connected the distribution of primes to the zeta
function through analytic continuation and a functional equation. Its
discussion of the zeros led to the critical-line hypothesis that now
bears his name
\cite[English translation, pp.~2--4]{riemann1859}. The present argument
concerns a zero-free half-plane for a family of $L$-functions; it does not
reach the critical line.

We work over $F=\mathbb Q(\sqrt{-3})$, with ring of integers $\OO$ and ideal
norm $\N\mathfrak a=|\OO/\mathfrak a|$. A finite-order Hecke character
$\eta$ is extended by zero to ideals not coprime to its modulus. For
$\Rea s>1$ its $L$-function is
\[
 L_F(s,\eta)=\prod_{\pp}
       \bigl(1-\eta(\pp)(\N\pp)^{-s}\bigr)^{-1}.
\]
We use its usual meromorphic continuation. Removing the factors at a fixed
finite set $S$ of prime ideals gives $L_F^S(s,\eta)$. The corresponding
finite factors are nonzero for $\Rea s>0$. As usual, a principal character
may have a pole at $s=1$.

\begin{theorem}\label{main}
All finite-order Hecke $L$-functions over $F$, and all
Dirichlet $L$-functions, are zero-free in the strict half-plane
\[
             \Rea s>\frac{3499999}{4000000}.
\]
The pole at $s=1$ for a principal character is allowed. The boundary is
$1/4000000$ smaller than $7/8$; no assertion is made about zeros on the
boundary line.
\end{theorem}

The deduction uses the results stated and proved unconditionally in
\cite{q30}; it requires no additional analytic conjecture.
Theorem~1.1 of that paper supplies the existing
$7/8$ bound. We also use Propositions~2.1, 11.3, and 15.2;
Lemmas~10.2, 11.1, 13.1, 14.3, 17.1, 17.6, 18.1, and 19.1; and the
coefficient, reflection, and zero-detection results of Sections~4--8.
When the perturbation changes a numerical range, we derive the needed
estimate from the general energy or moment bounds and the exact local
formulas in that paper. The foundational reflection, sieve, and moment
proofs are not reproduced here. Page references to \cite{q30} follow
its printed PDF pagination; references to its equations retain the
original numbering. The result here does not give a zero-free boundary
below $3/4$.

Let $Z$ tend to infinity. Write $\ell$ for the total length of the
compensation primes, and use $l_x,l_y$ to specify the column scales
$X=Z^{l_x}$ and $Y=Z^{l_y}$. These exponents have sum $M$, and the second exceeds the first
by $b$. The exponent $h$ determines the upper end of the central-row
range before a small buffer is added; $\st$ is the proposed zero-free
boundary. Increasing $\ell$ shortens both columns and shifts that upper
end. We fix the parameters as follows:

\begin{equation}\label{parameters}
\begin{gathered}
 \tau=10^{-6},\qquad \ell=\frac16+\tau,
 \qquad M=\frac56-\tau,\qquad b=\frac18,\\
 l_x=\frac{17}{48}-\frac\tau2,
 \qquad l_y=\frac{23}{48}-\frac\tau2,
 \qquad h=\frac{13}{16}+\frac{3\tau}{2},\\
 \st=\frac78-\frac\tau4=\frac{3499999}{4000000}.
\end{gathered}
\end{equation}
Thus
\[
 l_y-l_x=b,\qquad l_x+l_y=M,\qquad M+\ell=1,
 \qquad h=1-l_x+\ell.
\]
For the two column scales take $X=Z^{l_x}$ and $Y=Z^{l_y}$.
The adjustment leaves the affine exponent of the analytic signal
unchanged:
\begin{equation}\label{signal-exponent}
 C(s)=\frac{l_x}{2}+s-1+\frac h6=s-\frac{11}{16},
 \qquad C(\st)=\frac3{16}-\frac\tau4.
\end{equation}

A rightmost zero need not exist. We therefore measure the whole family
by the supremum
\[
 \bp=\sup\Bigl(\{1/2\}\cup
 \{\Rea\rho:L_F(\rho,\eta)=0,\ 1/2\le\Rea\rho\le1,
       \ \eta\text{ primitive of finite order}\}\Bigr).
\]
The supremum need not be attained. Throughout the proof we suppose, for
a contradiction, that $\bp>\st$, and put
$\Delta_\tau=\bp-\st>0$. The existing result gives
\begin{equation}\label{beta-range}
 \st<\bp\le\frac78,\qquad \kappa=\frac34,
 \qquad \bp\le\frac{1+\kappa}{2}.
\end{equation}
Thus the fixed value $\kappa=3/4$ is admissible in Lemma~18.1 of
\cite{q30}, whose uniform range is $\kappa\in[3/4,1]$. This choice uses
the existing bound, which is available before we perturb the parameters.

The proof compares two estimates for the same finite arithmetic probe.
The low-side estimate bounds its size by reflected-row energy and an
additive Gram bound. The rescaling coefficients absorb the cost of the
longer compensation primes. The high-side representation separates the
principal row from the remaining rows. After normalization, the principal
row gives the signal in \eqref{signal}, whose integrand contains
$H_\eta(s)/L_F^S(s,\eta)$. The reciprocal $L$-function is the link
between continuation of this signal and the location of zeros, as expressed
by Proposition~2.1 of \cite{q30}.

Sections~\ref{sec-count} and \ref{sec-contours} bound the nonprincipal
contribution. The witness capacities give a row count with an explicit
margin. A joint estimate for the numerator and error slots preserves that
margin after the full contour is moved. Section~\ref{sec-signal} extracts
the principal signal and bounds the outer rows. Section~\ref{sec-assembly}
then chooses the parameters without circular dependence and combines the
two estimates needed by the continuation principle to exclude the gap
$\Delta_\tau$.

\section{One finite probe and its Euler representation}\label{sec-probe}

We use the finite weighted arithmetic sum $I_\eta(X,Y,Z)$ defined in
(6.1) of \cite{q30}, with column scales $X,Y$ and large parameter $Z$.
Its arithmetic form gives the low-side bound. Its Euler representation
isolates the principal signal and allows us to estimate the other terms.
In the Euler representation, $u$ indexes the physical rows, and $u=1$
is the principal row.

A compensation-prime slot selects one prime from a prescribed window.
The modification combines divisibility conditions at the selected primes
with rescalings of the columns. Both estimates must apply to exactly this
finite sum, including its conjugate phases and coprimality restrictions.
We first define the modification as an arithmetic sum, then derive its
integral representation.

Fix a quotient $T$ of a ray class group before choosing a target character.
For each target $\eta$, use fixed arithmetic data $S,b_*,\xi$ allowed
by Sections~6 and 13 of \cite{q30}; here $b_*$ is a fixed ideal, and $S$
contains the relevant moduli and small primes. Also fix the nonnegative,
nonzero smooth annular weights $W_0,W_1$ and the Fourier normalization of
Section~6 of that paper. Denote its finite physical probe (6.1) by
$I_\eta(X,Y,Z)$. Inserting $\one_{D\mid cn^3}$ into its completed rows
gives $I_{\eta;D}$. The ideals $c,n$ retain the original coprimality
conditions and the characters retain their zero extension. These definitions specify the finite arithmetic expressions in (6.1).
The integral below is a representation of those expressions.

Choose a fixed even integer $K$ and split the compensation length equally:
$\ell_i=\ell/K$ and $P_i=Z^{\ell_i}$. Choose $K$ pairwise disjoint
compact supports in $(1,2)$, separated by positive gaps, and nonnegative,
nonzero smooth weights $W_i$ on them. The prime window is
\[
 \mathcal P_i(Z)=\{
  \pp\notin S:\pp\in1_T,\ W_i(\N\pp/P_i)\ne0\}.
\]
The condition $\pp\in1_T$ means that its ray class is the identity.
For $J\subseteq\{1,\ldots,K\}$ write
$p_J=\prod_{i\in J}\pp_i$, with the empty product equal to the unit ideal.
The modified probe is the same finite expression as (12.5) of
\cite{q30}, at the new scales:
\begin{equation}\label{finite-probe}
\begin{aligned}
 I_{\eta,\mathrm{mod}}(Z)
 ={}&\sum_{(\pp_i)\in\prod_i\mathcal P_i(Z)}
           \prod_i W_i(\N\pp_i/P_i)\\
 &\quad\times\sum_{J\subseteq\{1,\ldots,K\}}
      (-1)^{|J|}(\N p_J)^{-3/2}\overline{\eta(p_{J^c})}\\
 &\qquad\qquad\times I_{\eta;p_{J^c}}
       \bigl(X/\N p_J,Y/\N p_J,Z\N p_{J^c}\bigr).
\end{aligned}
\end{equation}
Each window retains its original scale $P_i$. The conjugate in
\eqref{finite-probe} pairs the marked contribution with the rescaled
contribution in the local compensation. Identity ray class makes the
characters in $\widehat T$ equal to $1$ at that prime; it does not make
$\eta(\pp)$ equal to $1$. The target may have any unit phase there.

For the Euler representation use the local data
$P_\pp,P_\pp^*,H_\pp,V,W,D$ of (7.11)--(7.12) in \cite{q30}, and denote
the complex variable called $x$ in its Section~16 by $s$. At $Q=\N\pp$
the complete local replacement is
\begin{equation}\label{local-replacement}
 G_\pp=
 \frac{
   (\overline{\eta(\pp)}Q^s-Q^{-w})(1-V)(1-W)P_\pp^*
       -Q^{-w}(1-VW)}{1-D}.
\end{equation}
Equation (16.4) of that paper reads
\[
 \overline{\eta(\pp)}Q^{s+z-1}P_\pp^*-Q^{z-w-1}P_\pp
 =\frac{1-D}{(1-V)(1-W)}Q^{z-1}G_\pp.
\]
Consequently the full correction is
\begin{equation}\label{full-correction}
\begin{aligned}
 H_{\eta,u,Z}(s,w,z)
 ={}&\sum_{(\pp_i)\in\prod_i\mathcal P_i(Z)}
       \prod_i\bigl[
        W_i(\N\pp_i/P_i)(\N\pp_i)^{z-1}G_{\pp_i}\bigr]\\
 &\quad\times
       \prod_{\substack{\pp\notin S\\
                  \pp\notin\{\pp_1,\ldots,\pp_K\}}}H_\pp.
\end{aligned}
\end{equation}
The local factors are evaluated at the current physical row $u$. Formula
\eqref{local-replacement} remains defined at zeros of $P_\pp$ or
$H_\pp$, so it also provides the continuation of the complete correction.
Disjoint supports ensure that each selected prime is replaced exactly
once, even though the nominal scales $P_i$ are equal.

To control the integral variables, we use the decay supplied by the
smooth weights. Let
$\widetilde W_0(|y|^2)$ be the plane Fourier transform of $W_0(|y|^2)$
for the self-dual measure of \cite{q30}, and set
\[
 \mathcal M(z)=\int_0^\infty\widetilde W_0(r)r^{z-1}\,dr,
 \qquad
 \widehat W_1(w)=\int_0^\infty W_1(r)r^{w-1}\,dr.
\]
The resulting weight is
\[
 \mathcal W(X,Y,Z;s,w,z)
   =X^{1/2-z}Z^{s+z-1}Y^{w-1}
      e^{(s+z-1)^2}\mathcal M(z)\widehat W_1(w).
\]
Equations (7.5), (7.13), and the finite replacements give the exact
identity
\begin{equation}\label{exact-high}
\begin{aligned}
 I_{\eta,\mathrm{mod}}(Z)
 ={}&\frac1{(2\pi i)^3}
   \int_{\Rea s=3}\int_{\Rea w=3}\int_{\Rea z=2}
       \mathcal W(X,Y,Z;s,w,z)\\
 &\quad\times\sum_u^{(6)}(\N u)^{-z}\overline{\xi(u)}
     \frac{\zeta_F^S(6z)L_F^S(w,\chi_\bullet(u))}
          {L_F^S(s,\eta\overline{\chi_\bullet(u)})}
       H_{\eta,u,Z}(s,w,z)\,dz\,dw\,ds.
\end{aligned}
\end{equation}
Here $(6)$ preserves the original physical restrictions of (7.5):
nonzero sixth-power-free elements, the chosen unit sector, and $(u,S)=1$.
The ideal character $\chi_\bullet(u)$ is the one in Lemma~4.1 of
\cite{q30}. Absolute convergence on the initial lines follows from the absolutely
convergent expansions of the original probe: only finitely many such
expansions are combined here. For each fixed $Z$, the slots, prime tuples, and
rescaling subsets are finite. The completed sum over $u$, the Euler
products, and the three Mellin integrals are infinite. The interchanges are justified by the absolute convergence used in
(7.5) of \cite{q30}. The fixed-tuple replacements form a finite
collection and preserve that convergence. Subsequent contour moves require the full-row
majorants proved in Sections~\ref{sec-contours} and \ref{sec-signal}. No local identity uses $\ell=1/6$. Thus
\eqref{exact-high} is still an exact high-side representation of
the finite probe \eqref{finite-probe} at the perturbed geometry.

\section{The cost of compensation on the low side}\label{sec-low}

Increasing the prime length changes the energy of the completed rows.
We therefore return to the general reflection estimate, whose dependence
on the lengths is explicit. Combining it with the additive Gram estimate
and the coefficient in \eqref{finite-probe} will show that the rescaling
absorbs the extra energy.

Fix a rescaling subset $J$ and write
\[
 d_J=\sum_{i\in J}\ell_i,\qquad
 \ell'=\ell-d_J,\qquad M'=M-2d_J,
\]
\[
 X'\asymp Z^{l_x-d_J},\qquad Y'\asymp Z^{l_y-d_J}.
\]
Then $0\le d_J\le\ell$ and $M'+\ell'-1=-3d_J$.
For a real number $t$, write $t_+=\max\{t,0\}$.

\begin{lemma}\label{row-energy}
The completed-row factor $B_{m,\sigma}^J$ in (15.1) of \cite{q30},
where $\sigma\in T$ is the finite ray class label, satisfies
\[
 \sum_{0<\N m\ll Q}|B_{m,\sigma}^J(Z)|^2
 \ll Z^{M'+(5\ell-1+d_J)_+/4+\eps},
 \qquad Q\asymp Z^{M'}.
\]
The constant may depend on the fixed windows and arithmetic data, and
is uniform in the moving rescaling prime tuple.
\end{lemma}

\begin{proof}
Apply the general reflected-energy estimate of Lemma~14.3 in
\cite{q30}. It decomposes the actual row into a powerful part, a part with fixed
support, and a remaining square-free part.
Let $O$ be the powerful-part length and $H$ the remaining-row length,
and put $\Delta_H=M'-O-H\ge-O(1/\log Z)$. The Gaussian subdivisions,
common-profile separation, and tail truncation on pp.~102--103 of that
paper leave the length identity unchanged. Denote their arbitrarily
small positive loss by $\eps_0$.

In the notation of Lemma~14.3 we have
$2A_0\le O+O(1/\log Z)$, $N_0\le A_0$, and $z_a\le\ell'$, with
\[
 T_{d_J}=2H+2A_0+2z_a-1-\ell'-\theta_N-N_0-3B_0.
\]
The identity $M'+\ell'-1=-3d_J$ gives
$T_{d_J}\le H-3d_J+O(\eps_0)$. The retained dual scales satisfy
$v+3\ell_b+e_\lambda\le T_{d_J}+\eps_0$, so
\[
                  \max(H,v+\ell_b)=H+O(\eps_0).
\]
The local parameter called $u$ in that energy lemma is
$\min\{v,z_a,(v+z_a)/3\}$; it is independent of the physical-row
notation in \eqref{exact-high}. If this minimum is $z_a$, the energy
is at most $M'-\Delta_H+O(\eps_0)$. Otherwise
$u\ge v/2-O(\eps_0)$ and
\[
 u+\ell_b+\frac{2e_\lambda}{3}
   +\frac{(T_{d_J}-v-3\ell_b-e_\lambda)_+}{2}
       \ge\frac{T_{d_J}}4-O(\eps_0).
\]
Substitute this in (14.14), using the nonnegative second numerator
in (15.3). In this second branch, the energy is bounded by
\[
 M'+\frac{1+3\ell'-2M'-2\Delta_H}{4}+O(\eps_0).
\]
Both bounds decrease, or stay constant, as $\Delta_H$ increases.
For an upper bound, we may therefore replace it by zero and absorb the
boundary error. Since
$1+3\ell'-2M'=5\ell-1+d_J$, the two branches give the asserted
positive-part bound. The common-profile kernel amplitude is extracted
before squaring. Powerful-row counting is already included in
Lemma~14.3, and its kernel saving is evaluated on the actual dual
support. Small losses, finitely many subsets, and logarithmically many
actual scales fit the stated $\eps$.

The change in lengths also affects the endpoint. When $d_J=\ell$,
$\ell'=z_a=0$ and
\[
 T_{d_J}\le1-6\ell+O(\eps_0)=-6\tau+O(\eps_0).
\]
A nonempty retained dual block has length exponent at least
$-O(1/\log Z)$. Choose the truncation loss sufficiently smaller than $\tau$, then
take $Z$ sufficiently large. The two length bounds are then incompatible,
so the retained block is empty. The discarded part is controlled by the uniform tail
estimate in Lemma~14.3, while the first branch also gives the claimed
bound. This handles the new zero-length endpoint directly.
\end{proof}

\begin{lemma}\label{low}
For every $\eps>0$,
\[
                 |I_{\eta,\mathrm{mod}}(Z)|
                    \ll Z^{C(\st)+\eps}.
\]
\end{lemma}

\begin{proof}
Set $P_a=Y'^2/Q=(\N b_*)^{-1}Z^b$. The scales for a fixed tuple
satisfy $Q,Y'\ge1$ and $P_a\ge1$: indeed
\[
 M-2\ell=\frac12-3\tau>0,\qquad
 l_y-\ell=\frac5{16}-\frac{3\tau}{2}>0.
\]
The additive Gram bound of Proposition~15.2 in \cite{q30} permits
these general polynomial scales. Moreover,
\[
 2b-l_y+d_J\le-\frac1{16}+\frac{3\tau}{2}<0.
\]
Hence its $P_a^2/Y'$ term is absorbed by the $P_a^{1/6}$ term, and
the squared Gram norm is at most
$Z^\eps(Q/Y')P_a^{1/6}$ times a fixed height polynomial.

Cauchy--Schwarz with Lemma~\ref{row-energy}, including the
$Q^{-1/2}$ factor in the original separation formula, gives the
fixed-tuple exponent
\[
 \frac{l_x-d_J}{2}+\frac b{12}
                   +\frac{(5\ell-1+d_J)_+}{8}.
\]
Summing over the rescaling tuples costs at most $Z^{d_J+\eps}$.
The coefficient in \eqref{finite-probe} contributes $Z^{-3d_J/2}$.
Combining these with the fixed-tuple estimate gives the exponent
\[
 L-d_J+\frac{(5\ell-1+d_J)_+}{8},
 \qquad L=\frac{l_x}{2}+\frac b{12}=C(\st).
\]
Because $\ell<1/5$, the function
$-d_J+(5\ell-1+d_J)_+/8$ starts at zero and is strictly decreasing;
its two slopes are $-1$ and $-7/8$. Every subset therefore has
exponent at most $L$. Rapid decay of the smooth weights absorbs the
fixed height polynomials, and the total small loss can be chosen below
the prescribed $\eps$.
\end{proof}

\section{Witness capacity and the central margin}\label{sec-count}

For the high-side estimate, we count rows by their norm, rather than by
the rescaling subsets used in the preceding section. Write
$\N u\asymp U=Z^d$, where $d$ is the row-norm exponent, and fix
$d_{\min}=1/100$. In a zero-detection box, $R$ denotes the ideal
row-count exponent. The bound $\#\mathcal B\ll_\eta U^{R+\eps_{\rm cnt}}$
must also account for the moment, witness, rounding, logarithmic, and
height losses. Those losses belong to the allowance $\eps_{\rm cnt}$,
which we keep separate from $R$.

The zero-detection argument in the source supplies inverse and ordinary
witnesses. Their moment estimates give two bounds for the row count, both
of which can be strengthened using the compensation primes already in
the probe. Each moment estimate permits a certain total prime-slot length;
we call this its capacity. A selection of slots must respect that capacity
and have phases compatible with the witnesses. We arrange the phases
first, then compare the two counts where their bounds cross.

As in Section~6.1 of \cite{q30}, put
$\Theta=\langle\eta,\widehat T\rangle$. Within each target, this finite
character group and every slot coefficient are fixed independently of
the moving row. The direction of the witnesses is handled in two
steps, following Section~19.2, pp.~182--183. First, if necessary,
conjugate the two witnesses and their smooth profiles so that their
common character is
\[
                 \psi(n)=\nu(n)\overline{\chi_n(u)},
                 \qquad \nu\in\Theta.
\]
Relative to this character the physical prime-slot coefficient is
\[
 \overline{\nu(\pp)}\one_{\pp\in1_T}
   =\frac1{|T|}\sum_{\vartheta\in\widehat T}
          (\overline\nu\vartheta)(\pp).
\]
It is a fixed combination of $\Theta$-characters, and the marked
moment permits this negative direction. For the positive-direction
ordinary moment in Lemma~18.1, conjugate the entire product of the
two ordinary witnesses and all selected physical slots. Its absolute
square is unchanged; the common character becomes
$\overline\nu(n)\chi_n(u)$ and the slot coefficient becomes
$\nu(\pp)\one_{\pp\in1_T}$. We conjugate the whole product, rather than change the phase of
an individual factor.

These two conjugations preserve the zero extension, the original prime
support, and the independence of the profiles. Since $\Theta$ is a group, membership
of the induced character in $\Theta$ is unchanged. Every retained
large row ramifies at a prime outside $S$, so its induced character
is outside $\Theta$, as Lemma~18.1 requires when the slot length is
positive. The witness heights may differ from the physical heights, so their
parameters require a separate bound.
After fixing the physical height and the selected slots, Lemma~4.5
handles the witness smooth parameters at the cost of a fixed height
power, uniform in the moving rows for the fixed arithmetic datum.
Indeed the physical normalized slot is
\[
 P_i^{-1/2}\sum_{\pp\in\mathcal P_i(Z)}
 \overline{\chi_\pp(u)}W_i(\N\pp/P_i)
                         (\N\pp/P_i)^{z-1}.
\]
Its annular profile is $W_i(y)y^{z-1}$ and its base-$U$ length is
$\ell_i/d$, since $U^{\ell_i/d}=P_i$. The character expansions above
identify exactly this polynomial with the prime factors in Lemmas~17.1
and 18.1. The disjoint supports are those fixed in
\eqref{finite-probe}, before the original zero masks are imposed.
Equal nominal lengths do not merge those supports. Thus the cited moment
inputs apply to the actual windows and profiles of this probe.

In a nonprincipal dynamic zero-detection box let $a\le\bp$ be its
real-part parameter, and put $\delta=2a-1$. Away from the bottom box,
\[
                 \frac1{50}<\delta\le\frac34;
\]
the bottom box uses $a=51/100$. Lemma~19.1 supplies an average
prime-amplitude parameter $q$ satisfying $0\le q\le\delta/2$.
Here $q$ is a dimensionless exponent, not a prime norm. For a
nonbottom box put $\alpha=5/6$ and $x=q/\delta\in[0,1/2]$.
The crossing of the two witness counts is described by
\begin{equation}\label{crossing}
\begin{aligned}
 D_x&=3-\frac{17x}{9},
 &P_x&=\left(2-\frac{8x}{9}\right)(1-x),\\
 J&=(\alpha-\delta)D_x+\delta P_x,\\
 R_*&=1-\delta+
        \frac{(\alpha-\delta)\delta P_x}{2J}.
\end{aligned}
\end{equation}
We have $D_x\ge37/18$, and $J$ has a positive lower bound on this
compact parameter range. The letter $J$ here is a real denominator;
it is unrelated to the rescaling subset in \eqref{finite-probe}.

\begin{lemma}\label{count}
Put $c_{\rm sup}=1/185$. Suppose that the available base-$U$ prime
length is at least $7/37+c_{\rm sup}$. Fix positive allowances
$\eps_m,\eps_w,\nu_0,\omega_{\rm slot}$ with
\[
 \eps_w\le10^{-3},\qquad
 \max\{\eps_w,\nu_0,\omega_{\rm slot}\}<c_{\rm sup}/16.
\]
Require $\omega_{\rm slot}$ to be no larger than the mesh supplied by
Lemma~18.1 of \cite{q30} for loss $\eps_m$ and the bounded length
ranges here, and require $\max_i\ell_i/d\le\omega_{\rm slot}$.
For each target with its fixed arithmetic data, final excluded set,
and fixed annular windows, there exist finite $A_{{\rm cnt},\eta}\ge0$,
an integer $L\ge0$, and a constant $C_\eta>0$ such that the following
raw bound holds in every retained nonbottom dynamic and amplitude bin
$\mathcal B$ under the detector's loss and height conditions:
\begin{equation}\label{raw-count}
 \#\mathcal B\le C_\eta(1+\log U)^L(1+H)^{A_{{\rm cnt},\eta}}
 U^{R_*+\eps_m+4\eps_w+6\nu_0+\omega_{\rm slot}}.
\end{equation}
Here $H\ge1$ bounds the retained physical and witness norm-twist
parameters, $\kappa=3/4$, and $1/50<\delta\le3/4$ as above.
The constants and orders may depend on the fixed arithmetic datum,
windows, slot count, strict moment margins, and prescribed losses,
but not on $U,H$, the moving rows or masks, the retained point, or any
later external tail order. They are fixed before a trial exponent in
$T_1=Z^\theta$ is chosen; the sufficiently large threshold may depend
on that exponent. Without selecting any slots, $t=1$ gives the same
type of raw bound with $R_*$ replaced by $1-2\delta/3$ and error
$\eps_m+4\eps_w$, without a prime-supply requirement.
\end{lemma}

\begin{proof}
The estimates below combine the crossing of the two short-witness
counts with the amplified no-slot bound. We choose the common
witness-length parameter to balance their exponents. For the selected
slots, we use the capacity calculation in the proof of Section~19.2.
Proposition~19.2 itself is specialized to the old positive gap, so its
statement cannot be applied directly here. The inverse witness uses Lemma~17.1, and the ordinary
witness uses Lemma~18.1; the latter's requirement
$\bp\le(1+\kappa)/2$ holds by \eqref{beta-range}.

Choose $t\in[1,3/2]$ as below, before selecting the witnesses.
Proposition~8.3 of \cite{q30} supplies a common presentation and actual
lengths $(r,m)$. Write $M_r$ for its normalized inverse witness and
$S_m$ for its normalized ordinary witness. Its saturation error constants
are absolute on these bounded ranges. Reduce its preliminary loss using those constants and
then increase the threshold to absorb the $O(1/\log U)$ errors. This
gives
\[
 \begin{gathered}
 r+m\ge t-\eps_w,\quad r\le t+\eps_w,\quad
 r\ge t-1/2-\eps_w,\quad 0\le m\le1/2+\eps_w,\\
 |M_r|^2\gg_\eta U^{\delta r-\eps_w},\qquad
 |S_m|^2\gg_\eta U^{\delta m-\eps_w}.
 \end{gathered}
\]
For example, if $C_*$ bounds all its saturation error constants, take
the preliminary loss at most $\eps_w/[2(C_*+1)]$; $C_*$ precedes the
target. The detector's literal height restrictions are verified in
Section~\ref{sec-assembly}.

Order the slots by decreasing nonnegative amplitude exponent, and
select only positive-exponent slots. If the request is $z$, a
fractional selection has weighted amplitude at least $qz$: if all
positive slots are exhausted, their full amplitude is $q\ell/d\ge qz$.
Discarding the last fractional slot loses at most
$\delta\omega_{\rm slot}$ in the squared exponent. Before rounding,
reduce a positive moment capacity $c$ to $z=c-\nu_0$; the total loss
relative to $2qc$ is at most
\[
       2q\nu_0+\delta\omega_{\rm slot}
                         \le\nu_0+\omega_{\rm slot}.
\]
The actual whole-slot length is at most $z$, even when there are too
few positive slots. Zero-exponent and error slots are used only through
upper bounds. Use the two direction changes above: first put witnesses
and physical slots in the common negative direction, then conjugate
the complete ordinary product. The respective coefficients are
$\overline\nu(\pp)\one_{\pp\in1_T}$ and
$\nu(\pp)\one_{\pp\in1_T}$, both independent of the moving row.
Zero extensions, independent windows, and the full-support
conditions are retained.

The ordinary-witness capacity is exactly
\[
           \frac{1-2m}{6\kappa}=\frac{2(1-2m)}9.
\]
At $\kappa=3/4$ it equals the baseline capacity used in the crossing
calculation of Section~19.2. There is no need for the $\Delta/4$
inflation introduced there by $\kappa=3/4+2\Delta$.
Here $Q$ denotes the product of the selected normalized prime
polynomials, in the common character direction just fixed.
The marked moment at effective row width one bounds
$\sum_u|M_r Q|^2$ by $U^{1+\eps_m}$ times its fixed height factor
when $r+2z\le1-c_1$ and $2r+8z\le3-c_2$.
The ordinary moment bounds $\sum_u|S_m^2Q|^2$ by the same power
when $2m+6\kappa z\le1$; its divisor is $|S_m|^4$.
At zero slot length its assertion holds for all the bounded $m$ here.

If $r\ge1$, the no-slot sixth-power amplified moment of Lemma~17.6,
equation (17.88), gives exponent
\[
 1-\alpha+(\alpha-\delta)r+\eps_m+\eps_w
 \le L(t)+\eps_m+(1+\alpha)\eps_w,
 \qquad L(t)=1-\delta+(\alpha-\delta)(t-1).
\]
Its form $e(r)=\max\{1,(1+5r)/6\}$ is uniform as $r$ approaches one.
If $m\ge1/2$, the ordinary no-slot count has exponent at most
$1-\delta+\eps_m+2\eps_w$. In the remaining range $r<1,m<1/2$,
the two ideal counts at their full capacities are
\[
 A_I(r)=1-\delta\{x+(1-x)r\},
\]
\[
 S_t(r)=1-\delta\{4x/9+(2-8x/9)(t-r)\}.
\]
Replacing $m$ by $t-r-\eps_w$ costs at most $2\eps_w$ because the
ordinary count decreases with $m$ and its derivative has absolute
value at most $2$. Its witness divisor costs another $2\eps_w$.
The crossing occurs at
\[
 r_*(t)=\frac{(2-8x/9)t-5x/9}{D_x},
 \qquad \frac{23}{37}\le r_*(t)\le1,
 \qquad \frac13\le t-r_*(t)\le\frac12.
\]
Use the ordinary count for $r\le r_*(t)$ and the inverse count for
$r\ge r_*(t)$. For the latter, if $(1-r)/2>\nu_0$, request
$z=(1-r)/2-\nu_0$. Its whole-slot length $z'$ satisfies
\[
 1-r-2z'\ge2\nu_0>0,\qquad
 3-2r-8z'=4(1-r-2z')+(2r-1)\ge\frac9{37}.
\]
Thus one may fix $c_1=2\nu_0$ and $c_2=9/37$ before the target.
For the ordinary request $z=2(1-2m)/9-\nu_0>0$, one has
$1-2m-6\kappa z'\ge(9/2)\nu_0$.
The unreduced inverse capacity is at most $7/37$, and the ordinary
capacity in its chosen branch is at most $2/27+4\eps_w/9$.
Both are strictly below the available supply with the fixed slack
specified in the statement. Decrement, witness error, and slot mesh
are each smaller than $c_{\rm sup}/16$; no limiting equality is used.

The common short-witness bound is
\[
 R_{\rm short}(t)=1-\delta+
                   \frac{\delta P_x}{D_x}(3/2-t).
\]
Choose
\[
             t=1+\frac{\delta P_x}{2J}\in[1,3/2],
\]
which makes $R_{\rm short}(t)=L(t)$ and yields exactly $R_*$ in
\eqref{crossing}. If the inverse capacity is at most $\nu_0$,
the no-slot ideal exponent is at most $1-\delta+2\delta\nu_0$;
if the ordinary capacity is at most $\nu_0$, it is at most
$1-\delta+(9/2)\delta\nu_0$. Since $R_*\ge1-\delta$, these
cases fit the common error
$\eps_m+4\eps_w+6\nu_0+\omega_{\rm slot}$.
For $q=0$ the selected-slot amplification factor is $1$, so the
argument requires no division by an amplitude. Empty bins contribute zero; a vanishing slot
cannot belong to a positive-amplitude class.

For completeness, without slots take $t=1$. For $r\le1$, the two
ideal exponents are $1-\delta r$ and $1-2\delta(1-r)$, whose
smaller value is at most $1-2\delta/3$, with crossing $r=2/3$.
For $r>1$, the amplified estimate is at most $1-\delta$ up to
the displayed witness errors. This proves the asserted no-slot form.

To pass from the ideal exponent to the raw count, we must also retain
the height and multiplicity factors in the moment estimates. Lemmas~17.1, 17.6, and 18.1 give finite seminorm and
height orders for the fixed datum, uniformly in their bounded real
ranges, masks, and moving supports. Their application includes the
finite derivative profiles required by the parameter Sobolev inequality
of Lemma~4.5. A logarithmic scale derivative inserts $yW'(y)$;
a normalized twist derivative inserts $\log y$ on the same annulus.
Sum over rows before the nonnegative Sobolev integrals. If a particular
fixed-parameter moment has height degree $J_j$ and $h_j$ height
coordinates, covering their box costs at most $(1+H)^{J_j+h_j}$.
Take $A_{{\rm cnt},\eta}=\max_j(J_j+h_j)$ over these actual calls,
their derivative profiles, and the no-slot calls. The finite presentation
labels and slot subsets enter $C_\eta$; witness dyads and logarithmic
scale coverings contribute $(1+\log U)^L$.
The stated moment theorems supply each of these finitely many orders,
so the maximum requires no extra analytic input. It may depend on $\eta$.
We fix its orders and constants before choosing the trial height exponent.
Auxiliary external tails in the detector are made smaller than a fixed
fraction of its spike after that exponent is chosen; their constants
only enlarge the lower threshold. They do not change these internal
moment orders. This proves \eqref{raw-count} with the asserted dependencies.
\end{proof}

The next step absorbs the logarithmic and height factors into powers
of $U$. Reserve positive $\eps_{\log}$ and
$\eps_{{\rm cnt},{\rm ht}}$ within the count allowance, with
\begin{equation}\label{count-loss-budget}
 \eps_m+4\eps_w+6\nu_0+\omega_{\rm slot}
       +\eps_{\log}+\eps_{{\rm cnt},{\rm ht}}<\eps_{\rm cnt}.
\end{equation}
The logarithm in \eqref{raw-count} is at most $U^{\eps_{\log}}$
for sufficiently large $U$. For $U=Z^d$, $d\ge d_{\min}$, set
$H=Z^{2\theta_\eta}$ and require
\begin{equation}\label{count-height-ceiling}
 0<\theta_\eta\le
 \frac{d_{\min}\eps_{{\rm cnt},{\rm ht}}}
      {4(1+A_{{\rm cnt},\eta})}.
\end{equation}
Then $(1+H)^{A_{{\rm cnt},\eta}}\ll_\eta
U^{\eps_{{\rm cnt},{\rm ht}}}$, so
$\#\mathcal B\ll_\eta U^{R_*+\eps_{\rm cnt}}$.
The same allocations cover the no-slot count, after taking the finite
maximum of its height orders, and the bottom bin has the elementary
bound $\#\mathcal B\ll_\eta U$. The actual choice of
$\theta_\eta$ is deferred to Section~\ref{sec-assembly}.

Proposition~\ref{central-joint} combines this row count with the
analytic factors in the high-side representation. After moving the full high-side
correction and only then subdividing its retained values, it gives
the joint estimate required by (10.13) of \cite{q30}. The resulting
general exponent of (10.15), relative to $C(7/8)$, is
\begin{equation}\label{high-exponent}
\begin{aligned}
 E_\tau(d)={}&a-\frac78
       +h\left(\frac{17}{50}-\frac16\right)
       -al_y-\frac\ell2+q\ell\\
      &\quad+d\left(R+\frac\delta2-\frac{17}{50}\right).
\end{aligned}
\end{equation}
Here $R$ is the ideal count exponent of physical rows in the chosen
box. The additional $\eps_{\rm cnt}$, already including count-height
absorption, contributes $d\eps_{\rm cnt}$ to the final $Z$ exponent;
it is not included a second time in the joint analytic height factor.
For the same $a,q,R,d$, its change from the old geometry is exactly
\begin{equation}\label{variation}
       E_\tau(d)-E_0(d)
         =\tau\left(\frac1{100}+\frac\delta4+q\right).
\end{equation}
We now exhibit a margin large enough to absorb this increase.

\begin{lemma}\label{endpoint}
For the adaptive witness count $R=R_*$, the endpoint $d=h$ satisfies
\[
 E_\tau(h)-\left(\bp-\frac78\right)
          \le-\frac{49}{440640}+\frac{15\tau}{8}<0.
\]
\end{lemma}

\begin{proof}
The complete square certificate (20.9) of \cite{q30} applies for
$0\le\delta\le\alpha$ and $0\le x\le1/2$. To display its positive
margin, set $y=1/2-x$, $v=51+41y$, and let $E_*$ denote the old
endpoint exponent. Its identity is
\begin{align*}
 10368vJ(-E_*)={}&(3+5y)\bigl\{(4v\delta-79)^2+49\bigr\}\\
 &+4y\bigl\{4v\delta[
       (1+3y)(15+32y)\delta+9-13y]+265+3485y\bigr\}.
\end{align*}
Every term is nonnegative, since $9-13y\ge5/2$. Using
$v\le17(3+5y)$ and $J\le5/2$ gives
\[
                         E_0(13/16)\le-\frac{49}{440640}.
\]
Increase $d$ from $13/16$ to $h=13/16+3\tau/2$, and then use
\eqref{variation}. This gives
\[
 E_\tau(h)\le-\frac{49}{440640}
       +\tau\left(\frac32R_*+\delta+q-\frac12\right).
\]
Since $J\ge\delta P_x$, we have
$R_*\le17/12-3\delta/2$. Together with $q\le\delta/2$ this bounds
the parentheses by $13/8$. Finally,
$\bp-7/8>-\tau/4$, so comparison with $C(\bp)$ costs at most a
further $\tau/4$.
\end{proof}

The other central ranges have larger fixed margins. The bottom box
uses $\delta_0=1/50$ and $R=1$, with old endpoint bound $-7/1200$;
the geometric change and comparison cost at most $32\tau/25$.
For $d\le1/2$ no selected primes are used: take
$R=76/75-2\delta/3$. Equation (20.11) gives $-49/14400$, and the
change and comparison cost at most $329\tau/400$. Thus put
\begin{equation}\label{central-margin}
                 m_A=\frac{49}{440640}-\frac{15\tau}{8}>10^{-4}.
\end{equation}
In the adaptive branch the slope in $d$ is positive and less than
$2$, so its endpoint controls $d\le h$. Choose
\[
              0<\zeta<\min\{1/96,m_A/32\}.
\]
Extending to $h+\zeta$ costs less than $2\zeta$. The same extension
still leaves enough compensation primes: since
$5\ell-h=1/48+7\tau/2$, on $1/2\le d\le h+\zeta$ we have
\[
                       \frac\ell d>\frac15>\frac7{37}.
\]
Since $1/5-7/37=2/185$, this implies the prime-supply condition
of Lemma~\ref{count} with $c_{\rm sup}=1/185$ on the entire range
where its selected slots are needed. The remaining fixed half of the
gap keeps every chosen decrement and rounding allowance strictly
inside the supply.

\section{Full contours and the joint central estimate}\label{sec-contours}

The contour lemmas must also be checked at the new boundary.
Lemmas~10.3--10.6 in \cite{q30} are stated only for boundary parameters
at least $7/8$. To extend the required arguments, we use the exact local
replacements and the full representation \eqref{exact-high}, checking
the ranges where the perturbation matters. The full correction is moved
before its retained values are subdivided by amplitude. This order lets
us use the row count without requiring the individual amplitude pieces
to be holomorphic.

Write $x_r=\Rea s$, $w_r=\Rea w$, and $z_r=\Rea z$. The amplitude
ratio $x=q/\delta$ of the previous section remains a separate real
parameter.

\begin{proposition}\label{extended-euler}
The complete correction \eqref{full-correction} has the analyticity
and full-height bounds required by Definition~10.1 of \cite{q30}
in its first Euler region and in the region $\mathcal D_{2,\tau}$
below. One final excluded set $S$ can make the principal slice satisfy
\[
            |H_\eta(s)-1|\le\frac12\qquad(\Rea s>\st).
\]
The central contours can still be moved for the full row sum, with
an external tail whose scale order is fixed independently of the
tail order chosen at the end.
\end{proposition}

\begin{proof}
The first Euler region is unchanged:
\[
 \mathcal D_1(\eps_0):\quad
 x_r\ge\frac{51}{100},\quad z_r\ge\frac{17}{50},\quad
 w_r\ge-\frac1{100},\quad x_r+w_r\ge1+\eps_0.
\]
For nonprincipal rows, move first on the global line
$x_r=\bp+20e$, and then, only within the buffered heights, to
\[
                      x_r=a+16e,\qquad w_r=1-a-6e,
\]
where $e>0$ is sufficiently small. Throughout the paths,
$x_r+w_r\ge1+10e$. The original denominator for a dynamic box is
zero-free in its buffered rectangle; global tails remain at
$x_r>\bp$ and use the global reciprocal bound. These paths need
only $\bp>51/100$, rather than $\bp\ge7/8$. The integral and
horizontal-edge tail estimates of Lemma~10.2, and the single
cumulative height allocation in (10.14), are unchanged. After moving the complete correction, we subdivide its values by
amplitude only to obtain pointwise upper bounds.

For the principal row enlarge the second Euler region to
\[
 \mathcal D_{2,\tau}:\quad
 x_r\ge\st,\qquad z_r\ge\frac{33}{200},\qquad
 w_r\ge\frac{19}{20}.
\]
This replaces the second region of Definition~10.1 and the
principal-slice contraction half-plane in (10.2). In the exact
local expressions (7.11), (7.17), the denominators $1-R,1-V,1-D$
stay uniformly nonzero: for $Q=\N\pp\ge2$,
\[
 |D|\le Q^{-\st},\qquad |V|\le Q^{-99/100},\qquad
 |R|\le Q^{-56/25+3\tau/2}.
\]
The largest off-row defect exponent and the largest dividing-prime
exponent are, respectively,
\[
 1-x_r-w_r-6z_r\le-\frac{363}{200}+\frac\tau4,
 \qquad
 1-x_r-w_r\le-\frac{33}{40}+\frac\tau4.
\]
For clarity, the ramified boundary terms $J_j$ in (7.17), for
$j=1,\ldots,5$, have exponents bounded above by
\[
 -\frac{73}{40}+\frac\tau4,\quad
 -\frac98+\frac{3\tau}4,\quad
 \left(-\frac{63}{40}+\frac{3\tau}4,-\frac32+\tau\right),\quad
 -\frac{39}{20}+\tau,\quad -\frac94+\frac{3\tau}2.
\]
These and the remaining off-row terms in (7.11), (7.17) are
strictly smaller than the corresponding largest exponents above.
Writing $\gamma_0=163/200-\tau/4>0$ and
$\gamma_1=33/40-\tau/4>0$, the positive local bounds are
$|H_\pp|\le1+O(Q^{-1-\gamma_0})$ off $u$ and
$|H_\pp|\le1+O(Q^{-\gamma_1})$ at primes dividing $u$.
The first product converges. For every $\eps>0$, the second is
$O_\eps((\N u)^\eps)$: outside a fixed finite set each factor is
at most $Q^\eps$, and the product of the distinct dividing primes
is at most $\N u$. These estimates are uniform in all imaginary
parts and unit phases. Thus the unselected Euler product is normally
convergent in the enlarged region with the stated full-height bound.
For $u=1$, enlarge the fixed cutoff in $S$ to make the product
uniformly close to $1$. The positive product-tail majorant is
$O(P_0^{-163/200+\tau/4})$, and gives
\[
                 \sup_{\Rea s>\st}|H_\eta(s)-1|\le\frac12.
\]
Further enlargement of $S$ uses the same positive tail bound,
without relying on cancellation between complex products.

The complete correction also has the full-row and full-height
bounds required by Definition~10.1. On the fixed real boxes where
(16.2)--(16.6) apply, $|G_\pp|\ll Q^B$ for a fixed $B$, with
height order zero. Bound the unselected Euler factors by the same
positive-product majorant and count the finite tuples. Then
\[
           |H_{\eta,u,Z}(s,w,z)|
                \ll_\eps(\N u)^\eps Z^{B'},
\]
where $B'$ is fixed before the final external tail order $N$.

The multiplier of the principal signal is the following slice:
\[
 H_\eta(s)=\prod_{\pp\notin S}
                   H_{\eta,1,\pp}(s,1,1/6).
\]
At $w=1,z=1/6$ the exact local factor gives
\[
 H_{\eta,\pp}(s)-1
   =O\bigl((\N\pp)^{-2}
        +(\N\pp)^{-1-\Rea s}
        +(\N\pp)^{3-6\Rea s}\bigr).
\]
The slice product therefore converges absolutely and uniformly
on compact subsets of $\Rea s>2/3$, and is holomorphic there; this
contains $\Rea s>\st$.

We next bound the full row sum, keeping the polynomial orders
independent of the tail order to be chosen later. First take
$\N u\le Z^{h+\zeta}$. There are $O(Z^{h+\zeta})$ elements in
the fixed unit sector. On each fixed real contour box the displayed
positive correction bound, Hecke strip growth, and the appropriate
reciprocal estimate give a polynomial majorant after this entire row
sum, of the form
\[
                 |F_Z(\vec{t})|
                 \le C_\eta Z^{B_\eta}(1+|\vec{t}|)^{J_\eta^*}.
\]
Here $F_Z$ includes the scale factors and arithmetic terms but not
the external smooth kernel. The principal poles are extracted as
residues; the bounds for its remaining contours stay away from those
poles. For the unbounded range $\N u>Z^{h+\zeta}$, the absolute
row majorant on $(2,2,z_\infty)$ and its summable dyadic series are
proved in Proposition~\ref{outer-rows}. They give another fixed
polynomial majorant. Taking the maximum over these finitely many types of estimates fixes
finite $B_\eta,J_\eta^*$ for all discarded integrals and joins before
we choose $N$.
A buffered reciprocal is used only on its allocated domain; other
tails stay on the global or absolute lines and use their global
reciprocal bounds. When an integration axis is extended, the buffered estimate is still
used only inside its rectangle.

For these same fixed $W_0,W_1$, the external kernel is
\[
 e^{-(t_s+t_z)^2}
       |\mathcal M(z_r+it_z)|\,|\widehat W_1(w_r+it_w)|.
\]
It is rapidly decreasing in all three variables: the map
$(t_s,t_z,t_w)\mapsto(t_s+t_z,t_z,t_w)$ is invertible, and
logarithmic integration by parts gives every decay order for the
two Mellin transforms. For $\mathcal M$, its underlying Fourier
transform is smooth at zero and Schwartz at infinity; all the real
$z$ intervals here are compact subsets of $(0,\infty)$. Thus the
boundary terms vanish. Annular support gives the same conclusion for
$\widehat W_1$. After the $w$ residue the identical argument applies
to the two remaining variables. Any external annular separation
profile $W$ is likewise fixed; its all-order Fourier estimates follow
by integration by parts, as in Lemma~4.5. Pure twists translate Mellin
arguments and append a fixed invertible block triangular height map.
Internal coefficient measures controlled only in a fixed weighted
$L^1$ norm remain integrated inside $F_Z$; no pointwise trace estimate
is inferred for such a measure.

For the truncations below, choose a fixed $c>0$ no larger than any
coordinate multiplier in (10.14). It depends only on the fixed
height-allocation matrix and is independent of $Z,T_1,N$.
With $k$ external coordinates, take kernel decay order greater than
$N+J_\eta^*+k$. Integrating the resulting positive polynomial
majorant outside $\max|t_j|\le cT_1$, or on a slice $t_j=\pm cT_1$,
gives, respectively, the integrated and horizontal-trace bounds
\[
                    O_{\eta,N}(Z^{B_\eta}T_1^{-N}).
\]
A bounded real join interval only multiplies this bound by its
length. This proves the tails used in the contour argument of
Lemma~10.2 and Section~10.3 in the enlarged regions. All absolute
interchanges use these full-row majorants. Increasing $N$ requires higher seminorms of the same external
functions. The polynomial orders $B_\eta,J_\eta^*$ and the internal
moment orders remain fixed.
Only the retained integral is subsequently subdivided pointwise.
\end{proof}

The row count must be combined with a bound for the analytic numerator
and all error slots together. The conductor saving belongs to the whole
numerator. Estimating each slot as if it had its own saving would use
that saving more than once.

\begin{proposition}\label{central-joint}
Put $z_0=17/50$ and $\delta=2a-1$. Consider central nonprincipal
rows in one retained box of Section~8 of \cite{q30}, at real parts
\[
                  x_r=a+16e,\quad w_r=1-a-6e,\quad z_r=z_0.
\]
Fix an error-slot set $I$ and an amplitude row set at a retained
point. Suppose the row count is at most $C_\eta U^{R+\eps_{\rm cnt}}$
after \eqref{count-loss-budget}--\eqref{count-height-ceiling}, and the mean
main-slot amplitude exponent is $q$, assigning exponent zero to
error slots. The corresponding piece of the full correction satisfies
\begin{equation}\label{joint-bound}
\begin{aligned}
 \sum_u\bigl|L_F^S(w,\chi_\bullet(u))
                    H^{(I)}_{\eta,u,Z}(s,w,z)\bigr|
 \ll_\eta{}&U^{R+\delta/2+\eps_{\rm cnt}+\eps_c}\\
 &\times Z^{\ell(z_0-1/2)+q\ell+\eps_c}
                      (1+T_1)^{A_\eta}.
\end{aligned}
\end{equation}
The central analytic loss $\eps_c$ excludes $\eps_{\rm cnt}$ and
can be made arbitrarily small after the slot system is fixed.
The finite height order $A_\eta$ here includes only the remaining
analytic factors; the count height has already been absorbed into
$\eps_{\rm cnt}$. It does not change with the external tail order. In particular the amplitude term in the general exponent
is $g=q\ell$.
\end{proposition}

\begin{proof}
Within the retained dynamic region, we can separate the local factor
into a main slot and an error slot as follows. Proposition~16.1 in \cite{q30} applies for $51/100\le a\le1$
without fixing total slot length to $1/6$. Its proof gives
\[
 H_\pp-1=O(Q^{-1-10e})\quad(\pp\nmid u),\qquad
 H_\pp-1=O(Q^{-10e})\quad(\pp\mid u).
\]
For fixed positive $e$, a sufficiently large final cutoff makes
these local factors nonzero. Write the uncompensated Euler product
as $\mathscr H_{\eta,u}=\prod_{\pp\notin S}H_\pp$, and define here
\[
 \mathfrak B_i=
   \sum_{\pp\in\mathcal P_i(Z)}
      W_i(\N\pp/P_i)(\N\pp)^{z-1}\frac{G_\pp}{H_\pp},
\]
\[
 Q_i=-\sum_{\pp\in\mathcal P_i(Z)}
        W_i(\N\pp/P_i)(\N\pp)^{z-1}
                         \overline{\chi_\pp(u)},
 \qquad D_i=\mathfrak B_i-Q_i.
\]
The conjugate in $Q_i$ is the one in Proposition~16.1, and agrees
with the negative physical direction used in Section~\ref{sec-count}.
For $\pp\nmid u$, the local cancellation (16.8) gives
\[
            G_\pp/H_\pp+\overline{\chi_\pp(u)}
                                     =O(Q^{-51/100}).
\]
Thus $D_i$ is precisely the error slot jointly estimated in (16.9).
Disjoint windows give the exact finite expansion
\begin{equation}\label{finite-expansion}
 H_{\eta,u,Z}
   =\mathscr H_{\eta,u}
       \sum_{I\subseteq\{1,\ldots,K\}}
          \prod_{i\in I}D_i\prod_{i\notin I}Q_i
   =\sum_I H^{(I)}_{\eta,u,Z}.
\end{equation}
This is (16.10). The quotients are valid because the local factors are nonzero in the
dynamic region. That nonvanishing was proved only there. Outside the
region, continuation uses the complete factor $G_\pp$ in
\eqref{local-replacement}.

The buffered estimates for the retained box and Lemma~4.10 give
the reflected primitive-numerator small-power bound required by
Proposition~16.1. Equation (16.9) then estimates the numerator
and every error slot simultaneously:
\begin{equation}\label{joint-errors}
\begin{aligned}
 \left|L_F^S(w,\chi_\bullet(u))\prod_{i\in I}D_i\right|
 \ll{}&U^{\delta/2+12e+\lambda_{\rm joint}}(1+T_1)^A\\
 &\times\prod_{i\in I}P_i^{z_0-1/2}.
\end{aligned}
\end{equation}
Here $\lambda_{\rm joint}>0$ is prescribed before the target,
and the preliminary small powers in that proof are reduced to fit it.
To track the factors in this estimate, put
$A_*=a-1/2+6e$. The primitive functional equation contributes
$q_f^{A_*}$ and a height polynomial, and the remaining deleted Euler
factors cost at most $U^{6e+\lambda}$ because $w_r\ge-6e$.
For the full set $J_0$ of distinct strict ramified labels, (16.11)
gives once
\[
 q_f\ll_S\N u\prod_{\pp\in J_0}Q^{-(j_\pp-1)}.
\]
Each such label has $j_\pp\ge2$ and
$-w_r-A_*=-1/2$. Hence its allocated local factor is at most
$Q^{z_0-1/2}$, simultaneously for all these labels; coprime errors
have the stronger exponent $z_0-51/100$. Divisor-many ramified
choices and the other primitive small powers are included in
$\lambda_{\rm joint}$. These two contributions account for $6e+6e=12e$. The estimate therefore
requires no additional positive loss proportional to the prime length.
The labels are distinct because the underlying windows are disjoint.
The conductor saving and the allowance for deleted factors have thus
been used once for the full product of error slots.

For the remaining slots the normalization is explicitly
\[
 \widetilde Q_i=P_i^{-1/2}
   \sum_{\pp\in\mathcal P_i(Z)}
       \overline{\chi_\pp(u)}W_i(\N\pp/P_i)
                          (\N\pp/P_i)^{z-1},
 \qquad Q_i=-P_i^{z-1/2}\widetilde Q_i.
\]
To form the amplitude boxes, write Lemma~19.1's bound as
$|\widetilde Q_i|\ll U^{\epsilon_1}P_i^{\delta/2+C_{\rm pr}e}$,
with $C_{\rm pr}$ fixed on the real ranges. After $K$ is fixed, put
$\ell_{\min}=\ell/K$ and $d_{\max}=h+\zeta$, and choose
\[
 C_{\rm pr}e+\frac{d_{\max}\epsilon_1}{\ell_{\min}}
                                             <\frac\vartheta2.
\]
Increasing the threshold absorbs the fixed constant into the other
half of the width, giving $|\widetilde Q_i|\le P_i^{\delta/2+\vartheta}$.
The capped subdivision of Section~19.1 then has
$0\le g_i\le\delta/2$, with $|\widetilde Q_i|\ge P_i^{g_i}$
when $g_i>0$. This prime-estimate cost is already included in
$\vartheta$. In an amplitude box take
$|\widetilde Q_i|\le P_i^{g_i+\vartheta}$,
using a zero-exponent upper bound for low-amplitude slots. Assign
$g_i=0$ to each error slot. Then
\[
                       \sum_i\ell_i g_i=q\ell,
                       \qquad 0\le q\le\delta/2,
\]
by Lemma~19.1. Multiply the main-slot bounds into
\eqref{joint-errors}, use $\mathscr H_{\eta,u}\ll U^\eps$, and sum
over the rows. For example, if
$\mathscr H_{\eta,u}\ll U^{\lambda_H}$, it suffices to choose
\[
 12e+\lambda_{\rm joint}+\lambda_H\le\eps_c,
                       \qquad \ell\vartheta\le\eps_c.
\]
This gives \eqref{joint-bound}, with $\eps_{\rm cnt}$ paid separately.
There are a fixed $2^K$ error subsets. Witness presentations and dyads
have already been included in the raw count's constant and logarithm;
only the remaining amplitude subdivisions and physical norm dyads
enter the subsequent multiplicity allowance.

Apply this estimate at each retained point of the already moved
complete integral. Its outside real exponent is
\[
 l_x(1/2-z_0)+a+z_0-1-al_y-dz_0+(16-6l_y)e.
\]
Adding \eqref{joint-bound}, subtracting $C(7/8)$, and using
$h=1-l_x+\ell$ gives \eqref{high-exponent} as the ideal exponent.
The additional real charges, including the retained reciprocal, are
listed in Section~\ref{sec-assembly}. All auxiliary heights obey the
single allocation (10.14). Neither an error piece
nor a pointwise row set is shifted on its own, and neither needs to
be holomorphic outside the retained region. This establishes the
use of (10.13)--(10.16) for the perturbed geometry.
\end{proof}

\section{The principal signal and the outer rows}\label{sec-signal}

The central estimate bounds the nonprincipal rows in the central range.
The principal row supplies the signal: taking its two scalar residues
produces the function whose continuation detects zeros. We first identify
the positive normalizing factor and the errors in this extraction.
We then estimate the small and large nonprincipal rows on contours
adapted to their sizes. These bounds complete the treatment of the
physical rows.

\begin{proposition}\label{principal-signal}
After fixing even $K$, the windows, and the final excluded set,
the principal-row double residue produces the signal
\eqref{signal}. Its normalizing factor $c_SA_T(Z)$ is positive,
with inverse bounded by an arbitrarily small power of $Z$.
The other principal-row integrals and the local approximation
retain the respective positive savings $l_y/20$, $h/600$, and
$\kappa_P$.
\end{proposition}

\begin{proof}
We first check that the positive tuple bound for the selected factors
still holds at the new lengths. On the principal rectangle the four error
powers in (16.8) of \cite{q30} are
\[
            -x_r,\qquad -6z_r,\qquad
            4-5x_r-6z_r,\qquad 1-w_r-6z_r,
\]
all strictly negative. The positive tuple upper bound for the
correction is therefore $Z^{\ell z_r+\eps}$ at every height.

At the principal double-residue slice $u=1,w=1,z=1/6$, define
\[
                    B_\pp(s)=
                    \frac{G_\pp(s,1,1/6)}{H_\pp(s,1,1/6)}.
\]
This slice lies in the second Euler region for the principal row,
where $H_\pp$ is nonzero. This quotient is used to normalize the principal row. For contour
movement in general, we continue to use the complete replacement \eqref{local-replacement}. The local identities
(16.12)--(16.13) give
\[
                 B_\pp(s)=-1+O((\N\pp)^{-\st}),
\]
since $5\st-3>\st$.

To normalize the prime sums, set
\[
 S_i(Z)=\sum_{\pp\in\mathcal P_i(Z)}
             W_i(\N\pp/P_i)(\N\pp)^{-5/6},
 \qquad
 A_T(Z)=(-1)^KZ^{-\ell/6}\prod_i S_i(Z).
\]
The fixed ray-class prime theorem, Lemma~13.1 of \cite{q30}, gives
\[
 S_i(Z)\sim\frac{c_iP_i^{1/6}}{|T|\log P_i},
 \qquad c_i=\int W_i(y)y^{-5/6}\,dy>0.
\]
Even $K$ makes $A_T(Z)>0$, and
$A_T(Z)\asymp(\log Z)^{-K}$. Once $K$ is fixed, choose
\[
                         \kappa_P=\frac{\st\ell}{2K}>0.
\]
The local approximation to the double residue has error
$O(Z^{-\kappa_P})$, uniformly on the entire global $s$ line.

Repeat the principal contour argument of Lemma~10.5 with the
extended region proved above. Keep $s_r=\bp+e$, first cross
$w=1$, and then, within its residue, cross $z=1/6$. No other
poles are crossed. Put
$r_S=\operatorname*{Res}_{v=1}\zeta_F^S(v)>0$. The product of the
two scalar residues is
\[
                  c_S=\widehat W_1(1)\mathcal M(1/6)r_S^2/6>0.
\]
The positivity of $\mathcal M(1/6)$ is the Fourier-pairing result
in (7.2), and the factor $1/6$ is the Jacobian in the residue of
$\zeta_F^S(6z)$. The unresidued $w$ integral saves $l_y/20$;
the remaining $z$ integral saves $h/600$.

Move the main term of the double residue to $s_r=2$. It yields
the height-independent signal
\begin{equation}\label{signal}
 f_\eta(Z)=\frac1{2\pi i}\int_{\Rea s=2}
       Z^{C(s)}e^{(s-5/6)^2}
                 \frac{H_\eta(s)}{L_F^S(s,\eta)}\,ds.
\end{equation}
We estimate the residue-approximation error on the original global
line, and move only the main term to the signal contour.
\end{proof}

\begin{proposition}\label{outer-rows}
After removing the principal row, the small rows
$0<\N u\le Z^{1/100}$ and the large rows
$\N u>Z^{h+\zeta}$ both retain a positive exponent saving.
Together with the central range they cover all physical rows.
\end{proposition}

\begin{proof}
For small rows $u\ne1$ use
\[
              x_r=\bp+e,\qquad w_r=1/2,\qquad z_r=17/50.
\]
These contour real parts lie in $\mathcal D_1(1/3)$ because $\bp>\st>5/6$;
the fixed margin $1/3$ thus remains valid in the new contradiction
range. At selected primes not dividing $u$, (16.8) gives
$G_\pp=O(1)$. At dividing primes its boundary exponents are
\[
 -1/2,\quad 3/2-2x_r,\quad
 (3/2-2x_r,2-3x_r),\quad 2-3x_r,\quad 3-5x_r.
\]
All are at most $1/2$, so $G_\pp\ll(\N\pp)^{1/2}$.
The divisor bound absorbs the factors associated with primes dividing
the row.
The original numerator and row-count bounds have combined power
$63/50<2$. The geometric exponent is
\[
 h\left(\frac{17}{50}-\frac16\right)-\frac{l_y}{2}
       =-\frac{79}{800}+\frac{51\tau}{100}.
\]
Allowing the conservative row-count cost $2/100$ and comparing
with $C(\bp)$ still leaves the saving
\[
                         \frac{63}{800}-\frac{19\tau}{25}>0.
\]

For large rows use the absolute lines
$(x_r,w_r,z_r)=(2,2,z_\infty)$, with fixed $z_\infty>2$.
The exact local formula gives $G_\pp=O_{z_\infty}(1)$ both off
and on the row. The positive unselected product is
$O_\eps((\N u)^\eps)$. Absolute ideal counting in each annular
slot, before any character cancellation, therefore bounds the full
selected-tuple sum by
\[
                 (\N u)^\eps\prod_iP_i^{z_\infty}
                       =(\N u)^\eps Z^{\ell z_\infty}.
\]
All three scalar Euler factors are absolutely bounded on these lines.
For $V=Z^{h+\zeta}$, dyadic lattice counting gives the convergent
series
\[
 \sum_{\N u>V}(\N u)^{-z_\infty+\eps}
 \ll\sum_{j\ge0}(2^jV)^{1-z_\infty+\eps}
 \ll V^{1-z_\infty+\eps},
 \qquad z_\infty>1+\eps.
\]
The convergent dyadic series bounds the entire infinite large-row
range, uniformly in the Mellin heights. The fixed smooth kernel is integrable,
so absolute convergence also justifies interchanging this series with
the integrals. Since $h=1-l_x+\ell$, the resulting exponent tends
to minus infinity with $z_\infty$. Its uniform upper bound is
\[
 \frac{53}{32}-\frac{3\tau}{4}
       +(h+\zeta)(1+\eps)-\zeta z_\infty+\eps.
\]
Fix $\zeta$ first and then a sufficiently large fixed
$z_\infty$. This gives a uniform positive saving for the large
rows.
\end{proof}

\section{Choosing the parameters and excluding the gap}\label{sec-assembly}

All estimates now refer to the same physical probe: the low-side bound,
the central-row estimates, the principal residue, and the bounds for the
outer rows. We can combine their positive margins once the parameters
have been chosen in the right order. The witness mesh must be fixed
before the number of prime slots. Increasing that number reduces the
saving in the principal approximation, so that saving must not determine
the earlier mesh.

First fix the geometry and keep the assumed family-wide contradiction
gap $\Delta_\tau$ fixed. The gap is fixed for the family, independently of the individual
target. Fix the central margin and $\zeta$. The chosen buffer gives
\[
 m_A-2\zeta>\frac{15m_A}{16}
        =10^{-4}+\frac{5861}{2350080000}.
\]
Divide the central loss allowance between moment-count losses
and later analytic subdivision losses, requiring their total to
be less than
\[
                              \frac{5861}{4700160000}.
\]
This preserves $\mu=10^{-4}$. Within the first allowance, choose
$\eps_{\rm cnt}$ and all six terms on the left of
\eqref{count-loss-budget}, with its strict inequality. We reserve room for the height and logarithmic factors now. Their
finite degrees will be determined after the target is fixed. Fix the witness losses,
$\nu_0$, and $\omega_{\rm slot}$ as in Lemma~\ref{count}.
The mesh in Lemma~18.1 depends only on the bounded real ranges and
$\eps_m$, not on the target or the later $K$. Then choose a sufficiently
large fixed even $K$ such that $2\ell/K<\omega_{\rm slot}$.
Throughout the adaptive row range, the available base-U prime length
stays above $1/5$, and
Lemma~\ref{count} controls both capacities. After fixing the
window weights the principal approximation saving is fixed at
\[
                               \kappa_P=\frac{\st\ell}{2K}>0.
\]

The moment losses chosen before $K$ are used in the central row count,
so they need only fit its allowance. They need not be smaller than
$\kappa_P$: the principal local approximation is a separate error and
uses no moment mesh.
Only after fixing $K$ do we shrink the dynamic-box widths,
amplitude-box widths, and local small powers. The central charges
must fit their remaining allowance; the low-side charges must
fit suitable portions of $\Delta_\tau/2$; the principal charges
must fit suitable portions of $l_y/20$, $h/600$, and $\kappa_P$.
More explicitly, with $\eps_{\rm pr}$ including the principal
small powers and its one normalization allowance, require
\[
 (1+h)e+\eps_{\rm pr}<l_y/80,\qquad
 e+\eps_{\rm pr}<h/2400,\qquad
 e+\eps_{\rm pr}<\kappa_P/4.
\]
Also take $e$ below the detector's bounds for the already prescribed
witness losses. These are choices after $K$; none asks the earlier
moment loss or mesh to be smaller than $\kappa_P$.
Detection boxes can be refined beyond the fixed moment mesh
without changing that mesh or $K$.

To combine the estimates, we must express all central losses in the
same base, changing from $U=Z^d$ to $Z$. Keeping $R$ as the ideal
count exponent, the total real-exponent loss after normalization is
\begin{equation}\label{central-loss-ledger}
 (16-6l_y)e+d\eps_{\rm cnt}+(1+d)\eps_c
                  +d\eps_d+\eps_p+\eps_{\rm norm}.
\end{equation}
Here $\eps_d$ is exactly the retained reciprocal allowance in
(10.16) of \cite{q30}: the buffered estimate for
$1/L_F^S(s,\eta\overline{\chi_\bullet(u)})$ is $O_\eta(U^{\eps_d})$.
This loss comes from the reciprocal factor. The row count contributes $d\eps_{\rm cnt}$;
the joint numerator, simultaneous conductor allocation, deleted
factors, unselected product, and main-slot amplitude upper bounds
contribute $d\eps_c+\eps_c$, as proved in
Proposition~\ref{central-joint}. The power $\eps_p$ covers the remaining pointwise or norm-dyadic
logarithmic multiplicities. The witness logarithm is already included in \eqref{count-loss-budget}.
Finally $\eps_{\rm norm}$ pays $(c_SA_T(Z))^{-1}\ll_\eta
Z^{\eps_{\rm norm}}$ once for this contribution. It is separate
from the primitive or deleted Euler factors. The fixed $2^K$
expansion and finite character labels cost constants.

If instead one writes $R_{\rm inc}=R+\eps_{\rm cnt}$ in
\eqref{high-exponent}, its $dR_{\rm inc}$ has already paid the
count and the separate $d\eps_{\rm cnt}$ in
\eqref{central-loss-ledger} must be omitted. We use the ideal
convention throughout. Likewise the count height already absorbed
in $\eps_{\rm cnt}$ is not charged again in a retained analytic
height power. The principal estimate has its own single normalization
allowance inside $\eps_{\rm pr}$; it does not incur the central
normalization charge a second time.
Since $d\le h+\zeta<1$, it suffices to arrange
\[
 (16-6l_y)e+\eps_{\rm cnt}+2\eps_c+\eps_d+
       \eps_p+\eps_{\rm norm}<\frac{5861}{4700160000}.
\]
The count share was reserved before $K$ and the remaining terms are
chosen after $K$ within the other share. Thus every raw and logarithmic loss has been included before we
choose the final analytic height.

Having retained every saving, take
\[
             m=\frac14\min\{\mu,l_y/20,h/600,\kappa_P\}>0.
\]
With the earlier moment losses and the later local small powers
chosen separately, we avoid circular dependence: increasing $K$ reduces
$\kappa_P$, which must not force another increase of $K$. All these
positive real parameters and finite structures, including the
large absolute line, are fixed before the target $\eta$.

For each target choose its fixed arithmetic data and enlarge
the excluded set to its final value. Complete every such
enlargement first, then use that same final $S$ for the physical
probe, Euler correction, $c_S$, $A_T$, and the two functions
\[
            J_\eta(Z)=\frac{I_{\eta,\mathrm{mod}}(Z)}{c_SA_T(Z)},
            \qquad f_\eta(Z)\text{ as in \eqref{signal}}.
\]
The low bound and logarithmic normalization give
\begin{equation}\label{normalized-low}
 |J_\eta(Z)|\ll_\eta Z^{C(\st)+\omega},
 \qquad \omega=\Delta_\tau/2\in(0,\Delta_\tau).
\end{equation}
Now fix all internal moment, seminorm, and Sobolev orders for this
datum. Lemma~\ref{count} gives $A_{{\rm cnt},\eta}$; let
$A_{{\rm ht},\eta}$ dominate the finitely many literal detector
height orders. The remaining retained analytic height degrees give
$A_\eta$, excluding the already paid count height. When a remaining
bound is expressed in $H=T_1^2$, include twice its degree in $A_\eta$.
The full-row majorants in Propositions~\ref{extended-euler} and
\ref{outer-rows}, with normalization and the remaining finite types,
give $B_\eta$. All these degrees are finite before a trial $\theta$
or the final external decay order is chosen.
Subject to the height ceilings below, combining the principal and all
nonprincipal estimates gives the same $m>0$ and, for every fixed $N$,

\begin{equation}\label{comparison-with-tail}
\begin{aligned}
 |J_\eta(Z)-f_\eta(Z)|\ll_{\eta,N}{}&
       Z^{C(\bp)-m}(1+T_1)^{A_\eta}\\
       &+Z^{B_\eta}T_1^{-N}.
\end{aligned}
\end{equation}
Here the analytic height is an auxiliary parameter in the error
estimate. The functions $J_\eta$ and $f_\eta$ themselves do not
depend on it.

Let $\eps_{\rm ht}>0$ be the minimum of the finitely many detector
height allowances fixed with the preliminary witness losses, and put
\[
 \theta_{0,\eta}=
   \frac{d_{\min}\eps_{\rm ht}}{20(1+A_{{\rm ht},\eta})}>0.
\]
For $T_1=Z^\theta$ with
$0<\theta\le\min\{d_{\min}/100,\theta_{0,\eta}\}$ and
$U\ge Z^{d_{\min}}$, sufficiently large $Z$ gives
\[
 T_1\le U^{1/100},\qquad
 (1+T_1)^{A_{{\rm ht},\eta}}
 \le2^{A_{{\rm ht},\eta}}Z^{d_{\min}\eps_{\rm ht}/20}
 \le U^{\eps_{\rm ht}/10}.
\]
These inequalities verify the height hypotheses of the pointwise
dyadic estimates. They also make the detector error
$U^{-189/100+o(1)}T_1^2$ tend to zero. All finitely many retained
physical and witness heights are bounded by a fixed multiple of
$T_1$, hence by $H=T_1^2=Z^{2\theta}$ after increasing the threshold.
The matrix allocation (10.14) bounds all added Mellin frequencies
simultaneously, with total at most $T_1/2$ in each buffered argument.
The same allowance applies throughout the later estimates.

Choose the actual height by
\[
 0<\theta_\eta\le\min\left\{
 \frac{d_{\min}}{100},\theta_{0,\eta},
 \frac{d_{\min}\eps_{{\rm cnt},{\rm ht}}}
      {4(1+A_{{\rm cnt},\eta})},
 \frac{m}{4(A_\eta+1)}\right\},
 \qquad T_1=Z^{\theta_\eta}.
\]
This both absorbs the count height in its preallocated allowance
and gives $(1+T_1)^{A_\eta}\ll_\eta Z^{m/4}$ for the remaining
retained factors. Only then choose a fixed integer $N_\eta$ with
\[
                   B_\eta-N_\eta\theta_\eta<C(\bp)-m/2.
\]
Auxiliary external orders in the detector and prime estimates may
now also be increased so their discarded parts fit strict power
gaps. Only the derivative and decay orders of the fixed smooth profiles
are increased. Their constants can be absorbed by the final threshold.
This does not require applying an internal moment estimate to a newly
differentiated external profile. Consequently neither $A_\eta$, $A_{{\rm cnt},\eta}$,
$A_{{\rm ht},\eta}$, nor $B_\eta$ is reset by this last choice.
Equation \eqref{comparison-with-tail} becomes
\begin{equation}\label{normalized-comparison}
               |J_\eta-f_\eta|\ll_\eta Z^{C(\bp)-m/2}.
\end{equation}
Throughout, $W,W_0,W_1$, the slot windows, and the functions
$J_\eta,f_\eta$ remain fixed. The later orders refer only to higher
derivative and decay estimates of the fixed external smooth profiles.
The outer slot and presentation types are finite; the internal
completed rows, Euler products, and Mellin integrals retain the
absolute convergence and majorants proved above.

The order of dependence can now be stated explicitly:
$m,\omega$ and all real losses are fixed before $\eta$; for each
target the finite degrees and constants precede $\theta_\eta$;
then come $N_\eta$ and the lower threshold. Proposition~2.1 of
\cite{q30} requires the positive savings $\omega$ and $m/2$ to be
independent of the target, but explicitly permits target-dependent
constants, excluded sets, and thresholds. Since neither function
contains the auxiliary height, it also permits this target-dependent
choice of $\theta_\eta$, precisely as in Lemma~11.1. The height degrees remain target-dependent throughout; their values
are never required to be uniform over the family.

We can now apply the common-signal continuation principle,
Proposition~2.1 of \cite{q30}, which allows $\sigma_0=\st$.
The exponent $C$ has slope one; the multiplier $H_\eta$ is
holomorphic and satisfies $|H_\eta-1|\le1/2$ on the new
half-plane; and the same $J_\eta$ satisfies both
\eqref{normalized-low} and \eqref{normalized-comparison}.
The first estimate supplies the low bound, and the second compares
the same normalized probe with the signal. Together they are the two
estimates required by the continuation principle. It excludes $\bp>\st$ by choosing a target
with an actual zero sufficiently close to the family supremum.
Attainment of the supremum is unnecessary. The assumed gap
therefore leads to a contradiction, and
\[
                                      \bp\le\st.
\]

\subsection*{From primitive Hecke characters to Dirichlet characters}

The first step of Proposition~11.3 of \cite{q30} extends the
Hecke conclusion to imprimitive characters. Passing from a
primitive character to an induced presentation changes only
finitely many factors
$1-\eta(\pp)(\N\pp)^{-s}$, which are nonzero for $\Rea s>0$.

To transfer the result to Dirichlet characters, we distinguish the
deleted rational primes from the prime ideals above them. Let $\chi$ be a primitive Dirichlet character of
conductor $q_\chi$. Let $T_{\mathbb Q}$ be the rational primes
dividing $3q_\chi$, and let $T_F$ consist of all prime ideals
above them. These sets are used only in this transfer identity.
On ideals coprime to $3q_\chi$ define
$\eta(\mathfrak a)=\chi(\N\mathfrak a)$, and extend by zero
elsewhere. If $\alpha\equiv1\pmod{q_\chi\OO}$, then
$\N\alpha\equiv1\pmod{q_\chi}$, so this defines a finite-order
ray class character.

Write $\chi_{-3}$ for the quadratic Dirichlet character of conductor $3$ associated with $F=\mathbb Q(\sqrt{-3})$.

For $p\notin T_{\mathbb Q}$, a split prime has
$\chi_{-3}(p)=1$ and two prime ideals of norm $p$. Its Hecke
factor is $(1-\chi(p)p^{-s})^{-2}$. An inert prime has
$\chi_{-3}(p)=-1$ and one prime ideal of norm $p^2$, giving
\[
 (1-\chi(p)^2p^{-2s})^{-1}
       =(1-\chi(p)p^{-s})^{-1}(1+\chi(p)p^{-s})^{-1}.
\]
Absolute convergence of the Euler products, followed by
meromorphic continuation, yields
\begin{equation}\label{dirichlet-transfer}
 L_F^{T_F}(s,\eta)
      =L^{T_{\mathbb Q}}(s,\chi)
         L^{T_{\mathbb Q}}(s,\chi\chi_{-3}).
\end{equation}
All deleted factors are nonzero for $\Rea s>0$. In
$0<\Rea s<1$, and at $s=1+it$ with $t\ne0$, both Dirichlet
factors have no poles, so a zero of either produces a Hecke zero.
For $\Rea s>1$ absolute Euler convergence gives nonvanishing.

At $s=1$, a possible pole cancellation can occur only when one
primitive factor is principal. The other primitive factor is then
$\chi_{-3}$, for which
\[
        L(1,\chi_{-3})=\int_0^1\frac{dx}{1+x+x^2}
                          =\frac{\pi}{3\sqrt3}>0.
\]
The additional finite Euler factors are also nonzero at $s=1$.
This rules out cancellation of a zero by a pole at that point.
The same argument with nonzero finite Euler factors handles
imprimitive Dirichlet characters. This proves
Theorem~\ref{main}, with the strict half-plane $\Rea s>\st$.

\begin{remark}[Further improvements]\label{prospect}
A substantially larger perturbation would require estimates that
support the altered lengths:
for example, stronger reflected-row energy or a sharper joint
numerator and error-slot bound, or greater usable capacity in the
inverse and ordinary witness moments. The Euler-region checks,
outer-row bounds, and principal approximation would also have to
remain compatible with that geometry. These estimates have not
been established here for such a change. The present fixed
perturbation provides no automatic iteration to the Riemann
Hypothesis and no guarantee that successive improvements would
reach the critical line. Broader arithmetic questions, including
Hilbert's twelfth problem, are longer-term aspirations for the
program; this argument claims no progress on that problem.
\end{remark}

\section*{AI Closure}

This work developed through human--AI dialogue. Human questions and
hints guided the choice of the small parameter perturbation and the
requirement that both estimates concern the same finite probe.
AI assistance helped develop and explain the argument. The AI for Math
plugin supported code generation for the computations and Lean formalization.
The mathematical content of the paper is the deduction above
from the stated literature inputs; the collaboration does not
enlarge those inputs or the theorem's scope.

The Lean formalization of the three terminal statements has
passed compilation with Lean~4.34.1 and verification with
Comparator's native core.

\begin{thebibliography}{9}

\bibitem{q30}
OpenAI.
\emph{The Quasi-Riemann Hypothesis: A Zero-Free Half-Plane
$\Rea s>7/8$}.
September 30, 2026, 199 pp.
\href{https://github.com/openai/math/blob/adc7f1241b42e322a6451854ab7e4b4c146bf78a/preprints/The-Quasi-Riemann-Hypothesis-September-30-2026/paper.pdf}{Official PDF},
at Git commit
\texttt{adc7f1241b42e322a6451854ab7e4b4c146bf78a}.
The main inputs here are Sections~2 and 10--20, especially
printed pp.~97--107 and 183--196, together with the coefficient,
reflection, and zero-detection results of Sections~4--8.

\bibitem{riemann1859}
Bernhard Riemann.
\emph{Ueber die Anzahl der Primzahlen unter einer gegebenen Groesse}.
Monatsberichte der Berliner Akademie (1859), 671--680.
\href{https://www.claymath.org/library/historical/riemann/english/EZeta.pdf}{English translation}
by David R. Wilkins, hosted by the Clay Mathematics Institute.
The historical discussion above uses pp.~2--3 for the definition,
continuation, and functional equation, and pp.~3--4 for the hypothesis.

\end{thebibliography}
\end{document}
